
% pas de section

\begin{frame}
    \frametitle{Boucles}

    Exemple~:
    \inputminted{ocaml}{exemples/0intro/somme_liste_boucle.ml}
    \pause

    Avantages~:
    \pause
    \begin{itemize}
        \item Plus facile à compiler ou à interpréter~;
        \item Moins de problèmes mémoires.
    \end{itemize}
\end{frame}

\begin{frame}
    \frametitle{Fonctions récursives}

    Exemple~:
    \inputminted{ocaml}{exemples/0intro/somme_liste_rec.ml}

    \pause

    \begin{itemize}
        \item Avantage~: parfois préférable en terme de lisibilité de code.
        \pause
        \item Problème~: peut entraîner des dépassements de pile.
    \end{itemize}
\end{frame}

\begin{frame}
    \frametitle{Fonctions récursives~: problème}

    \mintinline{bash}{[raphael@arch ~]$ ocaml boucle.ml} \\
    \mintinline{text}{1249999975000000} \\ % 50 millions premiers entiers
    \mintinline{bash}{[raphael@arch ~]$ ocaml fonction_recursive.ml} \\
    \color{red}\mintinline{text}{Stack overflow during evaluation (looping recursion?).}
\end{frame}


\begin{frame}
    \frametitle{Passage des fonctions récursives aux boucles}

    \begin{itemize}
        \pause
        \item Étape classique de l’optimisation lors de la compilation de
            langages fonctionnels
        \pause
        \item Comment ça marche en pratique~?
    \end{itemize}
\end{frame}


\begin{frame}
    \frametitle{Passage des fonctions récursives aux boucles}

    \plan
\end{frame}
